
\documentclass[11pt,a4paper]{article}
% Shared preamble for Wolf lecture note transcriptions
% Usage: % Shared preamble for Wolf lecture note transcriptions
% Usage: % Shared preamble for Wolf lecture note transcriptions
% Usage: \input{preamble} in each chapter file
% Include guard: safe to \input multiple times (e.g. from main_wolf.tex)
\ifdefined\wolfpreambleloaded\endinput\fi
\def\wolfpreambleloaded{}

\usepackage{amsmath,amssymb,amsthm}
% Match Wolf book-class layout: a4paper, ~345pt text width
\usepackage[a4paper, textwidth=345pt, textheight=550pt]{geometry}
\usepackage{hyperref}
\usepackage{mathtools}

% ── Notation (consistent across all chapters) ──
\newcommand{\mcM}{\mathcal{M}}       % matrix algebra M_d(C)
\newcommand{\mcB}{\mathcal{B}}       % bounded operators B(H)
\newcommand{\mcA}{\mathcal{A}}       % multiplicative domain
\newcommand{\mcH}{\mathcal{H}}       % Hilbert space H
\newcommand{\diag}{\operatorname{diag}}
\newcommand{\one}{\mathbb{1}}        % identity operator
\newcommand{\CC}{\mathbb{C}}
\newcommand{\RR}{\mathbb{R}}
\newcommand{\NN}{\mathbb{N}}
\newcommand{\ZZ}{\mathbb{Z}}
\newcommand{\tr}{\operatorname{tr}}
\newcommand{\spec}{\operatorname{spec}}
\newcommand{\rank}{\operatorname{rank}}
\newcommand{\supp}{\operatorname{supp}}
\newcommand{\range}{\operatorname{range}}
\newcommand{\spn}{\operatorname{span}}
\newcommand{\FT}{\mathcal{F}_T}      % fixed point set
\newcommand{\XT}{\mathcal{X}_T}      % peripheral eigenspace
\newcommand{\mcX}{\mathcal{X}}       % peripheral eigenspace (bare)
\newcommand{\id}{\mathrm{id}}        % identity map
\newcommand{\ket}[1]{|#1\rangle}
\newcommand{\bra}[1]{\langle#1|}
\newcommand{\braket}[2]{\langle#1|#2\rangle}
\newcommand{\ketbra}[2]{|#1\rangle\!\langle#2|}

% ── Helper: properly undefine a theorem environment for redefinition ──
% Usage: \UndefEnv{theorem} before \newtheorem{theorem}{...}
\makeatletter
\newcommand{\UndefEnv}[1]{%
  % Remove from section reset list (if registered there)
  \@removefromreset{#1}{section}%
  \expandafter\let\csname #1\endcsname\relax
  \expandafter\let\csname end#1\endcsname\relax
  \expandafter\let\csname c@#1\endcsname\relax
}
\makeatother

% ── Theorem environments ──
\theoremstyle{plain}
\newtheorem{theorem}{Theorem}[section]
\newtheorem{proposition}[theorem]{Proposition}
\newtheorem{corollary}[theorem]{Corollary}
\newtheorem{lemma}[theorem]{Lemma}
\theoremstyle{definition}
\newtheorem{definition}[theorem]{Definition}
\newtheorem{example}[theorem]{Example}
\newtheorem{problem}[theorem]{Problem}
\theoremstyle{remark}
\newtheorem{remark}[theorem]{Remark} in each chapter file
% Include guard: safe to \input multiple times (e.g. from main_wolf.tex)
\ifdefined\wolfpreambleloaded\endinput\fi
\def\wolfpreambleloaded{}

\usepackage{amsmath,amssymb,amsthm}
% Match Wolf book-class layout: a4paper, ~345pt text width
\usepackage[a4paper, textwidth=345pt, textheight=550pt]{geometry}
\usepackage{hyperref}
\usepackage{mathtools}

% ── Notation (consistent across all chapters) ──
\newcommand{\mcM}{\mathcal{M}}       % matrix algebra M_d(C)
\newcommand{\mcB}{\mathcal{B}}       % bounded operators B(H)
\newcommand{\mcA}{\mathcal{A}}       % multiplicative domain
\newcommand{\mcH}{\mathcal{H}}       % Hilbert space H
\newcommand{\diag}{\operatorname{diag}}
\newcommand{\one}{\mathbb{1}}        % identity operator
\newcommand{\CC}{\mathbb{C}}
\newcommand{\RR}{\mathbb{R}}
\newcommand{\NN}{\mathbb{N}}
\newcommand{\ZZ}{\mathbb{Z}}
\newcommand{\tr}{\operatorname{tr}}
\newcommand{\spec}{\operatorname{spec}}
\newcommand{\rank}{\operatorname{rank}}
\newcommand{\supp}{\operatorname{supp}}
\newcommand{\range}{\operatorname{range}}
\newcommand{\spn}{\operatorname{span}}
\newcommand{\FT}{\mathcal{F}_T}      % fixed point set
\newcommand{\XT}{\mathcal{X}_T}      % peripheral eigenspace
\newcommand{\mcX}{\mathcal{X}}       % peripheral eigenspace (bare)
\newcommand{\id}{\mathrm{id}}        % identity map
\newcommand{\ket}[1]{|#1\rangle}
\newcommand{\bra}[1]{\langle#1|}
\newcommand{\braket}[2]{\langle#1|#2\rangle}
\newcommand{\ketbra}[2]{|#1\rangle\!\langle#2|}

% ── Helper: properly undefine a theorem environment for redefinition ──
% Usage: \UndefEnv{theorem} before \newtheorem{theorem}{...}
\makeatletter
\newcommand{\UndefEnv}[1]{%
  % Remove from section reset list (if registered there)
  \@removefromreset{#1}{section}%
  \expandafter\let\csname #1\endcsname\relax
  \expandafter\let\csname end#1\endcsname\relax
  \expandafter\let\csname c@#1\endcsname\relax
}
\makeatother

% ── Theorem environments ──
\theoremstyle{plain}
\newtheorem{theorem}{Theorem}[section]
\newtheorem{proposition}[theorem]{Proposition}
\newtheorem{corollary}[theorem]{Corollary}
\newtheorem{lemma}[theorem]{Lemma}
\theoremstyle{definition}
\newtheorem{definition}[theorem]{Definition}
\newtheorem{example}[theorem]{Example}
\newtheorem{problem}[theorem]{Problem}
\theoremstyle{remark}
\newtheorem{remark}[theorem]{Remark} in each chapter file
% Include guard: safe to \input multiple times (e.g. from main_wolf.tex)
\ifdefined\wolfpreambleloaded\endinput\fi
\def\wolfpreambleloaded{}

\usepackage{amsmath,amssymb,amsthm}
% Match Wolf book-class layout: a4paper, ~345pt text width
\usepackage[a4paper, textwidth=345pt, textheight=550pt]{geometry}
\usepackage{hyperref}
\usepackage{mathtools}

% ── Notation (consistent across all chapters) ──
\newcommand{\mcM}{\mathcal{M}}       % matrix algebra M_d(C)
\newcommand{\mcB}{\mathcal{B}}       % bounded operators B(H)
\newcommand{\mcA}{\mathcal{A}}       % multiplicative domain
\newcommand{\mcH}{\mathcal{H}}       % Hilbert space H
\newcommand{\diag}{\operatorname{diag}}
\newcommand{\one}{\mathbb{1}}        % identity operator
\newcommand{\CC}{\mathbb{C}}
\newcommand{\RR}{\mathbb{R}}
\newcommand{\NN}{\mathbb{N}}
\newcommand{\ZZ}{\mathbb{Z}}
\newcommand{\tr}{\operatorname{tr}}
\newcommand{\spec}{\operatorname{spec}}
\newcommand{\rank}{\operatorname{rank}}
\newcommand{\supp}{\operatorname{supp}}
\newcommand{\range}{\operatorname{range}}
\newcommand{\spn}{\operatorname{span}}
\newcommand{\FT}{\mathcal{F}_T}      % fixed point set
\newcommand{\XT}{\mathcal{X}_T}      % peripheral eigenspace
\newcommand{\mcX}{\mathcal{X}}       % peripheral eigenspace (bare)
\newcommand{\id}{\mathrm{id}}        % identity map
\newcommand{\ket}[1]{|#1\rangle}
\newcommand{\bra}[1]{\langle#1|}
\newcommand{\braket}[2]{\langle#1|#2\rangle}
\newcommand{\ketbra}[2]{|#1\rangle\!\langle#2|}

% ── Helper: properly undefine a theorem environment for redefinition ──
% Usage: \UndefEnv{theorem} before \newtheorem{theorem}{...}
\makeatletter
\newcommand{\UndefEnv}[1]{%
  % Remove from section reset list (if registered there)
  \@removefromreset{#1}{section}%
  \expandafter\let\csname #1\endcsname\relax
  \expandafter\let\csname end#1\endcsname\relax
  \expandafter\let\csname c@#1\endcsname\relax
}
\makeatother

% ── Theorem environments ──
\theoremstyle{plain}
\newtheorem{theorem}{Theorem}[section]
\newtheorem{proposition}[theorem]{Proposition}
\newtheorem{corollary}[theorem]{Corollary}
\newtheorem{lemma}[theorem]{Lemma}
\theoremstyle{definition}
\newtheorem{definition}[theorem]{Definition}
\newtheorem{example}[theorem]{Example}
\newtheorem{problem}[theorem]{Problem}
\theoremstyle{remark}
\newtheorem{remark}[theorem]{Remark}

\numberwithin{equation}{section}

\title{Chapter 4: Convex Structure \\ {\small (Wolf Lecture Notes - Transcription)}}
\date{}

\begin{document}
\maketitle
\tableofcontents

% Make this ``Chapter 4'' in an article-class transcription:
\setcounter{section}{3}

\section{Convex Structure}

\subsection{Convex optimization and Lagrange duality}

Admittedly, a section about optimization techniques is a bit off topic. However,
convex optimization and, in particular, \emph{semidefinite programming}\footnote{Depending on what community one comes from it needs some time to get used to the terminology: conic, semidefinite or linear programs are just optimization problems of a certain type.}
turned out to be an extremely useful tool in the context of quantum information
theory, so we will at least sketch the basic framework here.

The utility of these methods is threefold: (i) they provide powerful algorithms
to solve problems numerically. In contrast to many other numerical techniques
these algorithms are certifiable. That is, they return a result together with a
(usually negligeable) error bound which is essentially enabled by duality theory.
(ii) Duality theory allows to convert numerical results into analytic proofs and
often enough gives some analytic insight (if not a solution) without exploiting
any silicon brain. (iii) Having expressed a problem in terms of e.g.\ a
semidefinite program may be considered as a solution of the problem in very
much the same way as we consider a reduction to an eigenvalue problem to be a
solution. Moreover, it enables a classification of the problems complexity in
the sense of complexity theory since interior point methods provide a general
polynomial-time solution.

\noindent\textbf{Conic programs}
Consider finite-dimensional real Hilbert spaces $V$ and $V^\prime$. These may
be bare $\RR^d$'s but could also be spaces of Hermitian matrices or spaces of
polynomials of a certain degree. Suppose $K\subset V$ is a closed convex cone
which is pointed $(K\cap (-K)=\{0\})$ and has a non-empty interior. For any
linear map $T:V\to V^\prime$ and $c\in V$, $b\in V^\prime$ the optimization
problem
\begin{equation}\label{eq:conic-primal}
    \inf_{x\in K}\left\{\langle c|x\rangle \Bigm| T(x)=b\right\} \eqqcolon C_p
\end{equation}
is called a \emph{convex conic program}. Assigned to this is a \emph{dual
    problem}\footnote{The relation between the \emph{primal} problem in
    \eqref{eq:conic-primal} and the dual problem \eqref{eq:conic-dual} is more
    symmetric than it appears at first sight. In particular, the dual problem is a
    convex conic program in its own right. To see the symmetric structure note that
    $T(x)=b$ defines an affine set which can be parametrized by
    $\tilde{y}\in \RR^n$ via $x=e_0+\sum_{i=1}^n \tilde{y}_i e_i$ with suitable
    $e_i\in V$. Defining further $\tilde{b}_i \coloneqq \langle c|e_i\rangle$ we can
    rewrite the primal problem without equality constraints in a form similar to
    that of Eq.~\eqref{eq:conic-dual} up to a constant offset (since
    $\langle \tilde{b}|\tilde{y}\rangle=\langle c|x\rangle-\langle c|e_0\rangle$)
    and a sign which has to be spent to turn the inf into a sup.}
which reads
\begin{equation}\label{eq:conic-dual}
    \sup_{y\in V^\prime}\left\{\langle b|y\rangle \Bigm| c-T^{*}(y)\in K^{*}\right\}
    \eqqcolon C_d,
\end{equation}
where $T^{*}$ is the dual map (i.e., $\langle T^{*}(y)|x\rangle=\langle y|T(x)\rangle$)
and $K^{*}$ is the dual cone defined by
\[
    K^{*} \coloneqq \{x\in V \mid \forall z\in K:\ \langle z|x\rangle \ge 0\}.
\]
It holds in general that $C_p\ge C_d$ -- this is called \emph{weak duality}.
Strong duality, meaning $C_p=C_d$, holds if $T(x)=b$ is fulfilled for some $x$
which is in the interior of $K$. In this case the problem in
Eq.~\eqref{eq:conic-primal} is called strictly feasible and the supremum $C_d$
is attained in the dual. Similarly, strong duality holds if the dual problem in
Eq.~\eqref{eq:conic-dual} is strictly feasible (i.e., $c-T^{*}(y)$ is in the
interior of $K^{*}$ for some $y$). In this case the infimum $C_p$ is attained
in Eq.~\eqref{eq:conic-primal}.

For many conic programs there exist polynomial-time algorithms. Roughly
speaking, the requirements for that are that membership in the cone can be
decided efficiently (like for the cone of positive semidefinite matrices) and
that bounds on the norm of candidate solutions can be given.

\noindent\textbf{Semidefinite programs}
are special instances of convex conic programs. Here $T$ is a linear map
between spaces of Hermitian matrices and $K=K^{*}$ is the cone of positive
semidefinite matrices. Traditionally, semidefinite programs are phrased in a
slightly different way than in Eqs.~\eqref{eq:conic-primal} and
\eqref{eq:conic-dual}. The weak duality inequality is usually stated as
\begin{equation}\label{eq:sdp-weak-duality}
    \begin{aligned}
        \inf_{X\ge 0}\left\{\tr[F_0X] \Bigm| \tr[F_iX]=b_i \ (i=1,\dots,n)\right\}
        \ge
        \sup_{y\in \RR^{n}}\left\{\langle b|y\rangle \Bigm| F_0 \ge \sum_{i=1}^{n} y_i F_i\right\},
    \end{aligned}
\end{equation}
where $b\in \RR^{n}$, $F_0,F_1,\ldots,F_n$ as well as $X$ are Hermitian matrices
and ``$\ge$'' for matrices denotes the standard matrix ordering induced by the
cone of positive semidefinite matrices. Strong duality, meaning equality in
Eq.~\eqref{eq:sdp-weak-duality}, holds if at least one of the two optimization
problems is strictly feasible. That is, if there is either an $X>0$ with
$\tr[F_iX]=b_i$ $(i=1,\dots,n)$ or a $y$ such that $F_0>\sum_i y_iF_i$. As
mentioned in the previous paragraph, if one of the problems is strictly
feasible, then the extremum in the dual problem is attained.

If the extrema in Eq.~\eqref{eq:sdp-weak-duality} are both attained and equal
(i.e., in particular strong duality holds) then a simple condition called
complementary slackness relates optimal solutions $X^0$ and $y^0$:
\begin{equation}\label{eq:complementary-slackness}
    \left(F_0 - \sum_i y_i^{0} F_i\right) X^{0} = 0.
\end{equation}
That is, $X^0$ and $(F_0-\sum_i y_i^{0}F_i)$ have to have orthogonal
supports. Moreover, any $y^0$ is then optimal iff there exists an $X^0\ge 0$
such that Eq.~\eqref{eq:complementary-slackness} holds and the constraints
$\tr[F_iX^0]=b_i$ and $F_0\ge \sum_i y_i^{0}F_i$ are satisfied.

\subsection{Literature}

Convex optimization is described in detail in the textbook [?]. [?] is a good
source for more on semidefinite programs and more on conic programs, in
particular about complexity related questions, can be found in [?].

\end{document}
